\documentclass[10pt,a4paper]{amsart}
\usepackage[margin=19mm]{geometry}
\usepackage{amsmath,amssymb,mathtools}
\usepackage[scr=rsfs,cal=euler]{mathalfa}
\usepackage{xfrac}
\usepackage[unicode,hidelinks,bookmarks=false]{hyperref}
\newcommand{\ie}{i.e.\ }
\newcommand{\cf}{cf.\ }
\newcommand{\resp}{respectively\ }
\newcommand{\Subsection}{\subsection}
\providecommand{\nobreakdash}{\nobreak\hbox{-}}
\setlength{\emergencystretch}{2em}
\multlinegap0pt
\usepackage{bm}
\usepackage{bbm}
\usepackage{dsfont}
\usepackage{xspace}
\usepackage{cancel}
\usepackage{mathtools}
\usepackage{amsthm}
\makeatletter
\@ifundefined{theorem}{\newtheorem{theorem}{Theorem}}{}
\@ifundefined{corollary}{\newtheorem{corollary}[theorem]{Corollary}}{}
\makeatother
\def\R{\mathbb{R}}
\def\Sc{\mbox{Sc}\,}
\def\div{\mbox{div}\,}
\title[Calder\'on's inverse problem via Vekua theory]{Calder\'on's inverse problem via Vekua theory}
\author{Briceyda B. Delgado}
\date{Source: arXiv:2601.17313v1 (CC BY 4.0)}
\begin{document}
\maketitle

\noindent\emph{Source and licence.} Calder\'on's inverse problem via Vekua theory. Version arXiv:2601.17313v1 (CC BY 4.0), distributed under the Creative Commons Attribution 4.0 International licence, as recorded in the arXiv licence metadata for this exact version and shown on the abstract page https://arxiv.org/abs/2601.17313v1. Full rights notice: licenses/arxiv-2601.17313-CC-BY-4.0.txt.\par\smallskip

\noindent\textit{Editorial scope.} Counted unit: the Corollary whose source label is \texttt{cor:Cauchy-fundamental} in the article (source0.tex lines 390--401, source label \texttt{cor:Cauchy-fundamental}), delivered together with the article's own complete original proof of it (source0.tex lines 402--419, one \texttt{proof} environment of 18 lines). No mathematics is written, repaired or re-derived: statement and proof are the article's own text line for line. Required context retained uncounted: eq:Vekua\_equation (lines 142--144); constraint-1 (lines 309--312); constraint-2 (lines 309--312); th:Cauchy-theorem (lines 314--340); eq:scalar-Vekua (lines 370--377). Every definition, display and label of those lines is retained unchanged. Omitted from this package and not redistributed: 1--141, 145--308, 313--313, 341--369, 378--389, 420--1255. Nothing inside a retained excerpt is omitted or altered: every retained body is delivered exactly as the article's lines are written, so no omission receipt of any kind accompanies this package. Citations: the retained excerpts carry no citation command at all, so no bibliography entry of the article is transcribed and no reference is redistributed. Reference adaptation: none was needed and none was made; every reference inside the delivered unit resolves inside it, so no unresolved reference or citation is produced. Printed slips kept literal: no printed word, symbol, hypothesis or number of the counted statement or of its proof was altered. Layout and class: the article's own class is \texttt{article} with packages including amsmath, amsfonts, amssymb, amsthm, amstext, color, graphicx, hyperref; the wrapper is the lane's pinned \texttt{amsart} editorial wrapper at 10pt on A4, which restates the theorem-like environments the retained text needs and adds the article's own macro definitions that the retained text needs (\texttt{\textbackslash R}, \texttt{\textbackslash Sc}, \texttt{\textbackslash div}), with extra wrapper packages (bm, bbm, dsfont, xspace, cancel, mathtools). The delivered numbering is the unit's own. Assets: no retained span contains a figure environment, a graphics-inclusion command, a PDF-page inclusion, a table or a TikZ picture; no image file is copied into this package, the package declares no asset dependency and no image file is redistributed with it. Licence: version arXiv:2601.17313v1 of this article is distributed under the Creative Commons Attribution 4.0 International licence, as recorded in the arXiv licence metadata for this exact version and shown on the abstract page https://arxiv.org/abs/2601.17313v1. A full rights notice is delivered at licenses/arxiv-2601.17313-CC-BY-4.0.txt. Characters: every retained excerpt is pure ASCII, so the delivered engine, XeLaTeX with its default Latin Modern text font, reports no missing character.
\begin{align}\label{eq:Vekua_equation}
   D w=\alpha\, \overline{w}, \quad \text{ in }\Omega,
\end{align}

\begin{align}\label{constraint-1}
   \lim_{\epsilon\rightarrow 0} \int_{\partial B(x;\epsilon)} \Phi(y-x)\eta(y)v(y)\, ds_y=v(x),\\
   \lim_{\epsilon\rightarrow 0} \int_{\partial B(x;\epsilon)} \Psi(y-x)\eta(y)v(y)\, ds_y=v(x).\label{constraint-2}
\end{align}

\begin{theorem}\label{th:Cauchy-theorem}
 Let $\Omega\subset \mathbb{R}^n$ be a bounded domain with Lipschitz boundary. Suppose that $ u, v\in C^1(\overline{\Omega})$, and $\alpha\in L^{\infty}(\Omega)$. Then the following formulas are fulfilled
\begin{align}
   \text{Sc }\int_{\partial\Omega} u\, ds_y^* \, v&=\text{Sc }\int_{\Omega} u [ (D-\alpha C)v]-
   [(D-M^{\alpha}C)\overline{u}]\, \overline{v} \, dy\label{eq:integral-formula-1}\\
   \text{Sc }\int_{\partial\Omega} u\, ds_y^* \, v&=\text{Sc }\int_{\Omega} u [(D-M^{\alpha}C)v]-
   [(D-\alpha C)\overline{u}]\, \overline{v}\, dy.\label{eq:integral-formula-2}
\end{align}
Moreover, if $-\overline{\Phi}$ is a fundamental solution of the operator $D-M^{\alpha}C$ satisfying \eqref{constraint-1}, then we can recover the scalar part of $v$, $v_0$, as follows
\begin{align}\label{eq:scalar-BP}
\Sc\, \int_{\partial\Omega}\Phi(y-x) \eta(y) v(y)\, ds_{y}-\Sc\, \int_{\Omega} \Phi(y-x) [ (D-\alpha C)v] \, dy=\left\{
   \begin{array}{ll}
      v_0(x), \ & x\in \Omega,\\
      0,           & x\in \R^n\setminus \overline{\Omega}.
   \end{array}
\right. 
\end{align}
Analogously,  if $-\overline{\Psi}$ is a fundamental solution of the operator $D-\alpha C$ satisfying \eqref{constraint-2}, then 
\begin{align}\label{eq:scalar-BP-2}
  \Sc\, \int_{\partial\Omega}\Psi(y-x) \eta(y) v(y)\, ds_{y}-\Sc\, \int_{\Omega} \Psi(y-x) [ (D-M^{\alpha} C)v] \, dy=\left\{
   \begin{array}{ll}
      v_0(x), \ & x\in \Omega,\\
      0,           & x\in \R^n\setminus \overline{\Omega}.
   \end{array}
\right. 
\end{align}
\end{theorem}

\begin{align}\label{eq:scalar-Vekua}
  \text{Sc }\int_{\partial\Omega} \Phi(y-x)\, \eta(y) \, w(y) \,ds_{y}=\left\{
   \begin{array}{ll}
      w_0(x), \ & x\in \Omega,\\
      0,           & x\in \R^n\setminus \overline{\Omega},
   \end{array}
\right.
\end{align}

\begin{corollary}\label{cor:Cauchy-fundamental}
Under the same hypothesis than Theorem \ref{th:Cauchy-theorem}. If $h_0$ is a scalar fundamental solution of $-\Delta+|\alpha|^2+\div \vec \alpha$, with $D\alpha=-\div \vec \alpha$, then
\begin{align*}
  \text{Sc }\int_{\partial\Omega} (\nabla+\overline{\alpha})h_0(y-x)\, \eta(y) \, w(y) \,ds_{y}=\left\{
   \begin{array}{ll}
      w_0(x), \ & x\in \Omega,\\
      0,           & x\in \R^n\setminus \overline{\Omega},
   \end{array}
\right.
\end{align*}
for any $w$ solution of the Vekua equation \eqref{eq:Vekua_equation} whose boundary values are well-defined.
\end{corollary}

\begin{proof}
First, we will verify the following factorization
\begin{align}\label{eq:factorization-2}
    \left(-\Delta+|\alpha|^2-D\alpha\right)h_0=(D-M^{\alpha}C)(D-\alpha C)h_0.
\end{align}
Indeed,
\begin{align*}
(D-M^{\alpha}C)(D-\alpha C)h_0&=(D-M^{\alpha}C)(Dh_0-\alpha h_0)=D^2 h_0-D(\alpha h_0)-\overline{Dh_0}\alpha+\overline{\alpha h_0}\alpha\\
&=-\Delta h_0-D(\alpha) h_0- \nabla h_0\alpha+\nabla h_0\alpha+|\alpha|^2h_0.
\end{align*}
%\begin{align*}
%(D-\alpha C)(D-M^{\alpha}C)h_0&=(D-\alpha C)(Dh_0-\alpha h_0)=D^2 h_0-D(\alpha h_0)-\alpha\overline{Dh_0}+\alpha\overline{\alpha h_0}\\
%&=-\Delta h_0-D(\alpha) h_0- \nabla h_0\alpha+\nabla h_0\alpha+|\alpha|^2h_0.
%\end{align*}

By \eqref{eq:factorization-2}, if $h_0$ is a fundamental solution of $-\Delta+|\alpha|^2-D \alpha=-\Delta+|\alpha|^2+\div \vec \alpha$, then $\overline{\Phi}=(D-\alpha C)h_0$ will be a fundamental solution of $D-M^{\alpha}C$.
The result comes directly from \eqref{eq:factorization-2} and the Cauchy-type integral formula provided in \eqref{eq:scalar-Vekua}.
\end{proof}

\end{document}
